% from ega1-7-fr.tex
\oldpage[I]{159}
\begin{proposition}[7.2.2]
\label{I.7.2.2-fr}
Soient $X$, $Y$ deux $S$-préschémas, tels que $X$ soit réduit et $Y$ séparé
sur $S$. Soit $f$ une $S$-application rationnelle de $X$ dans $Y$, $U_0$ son
domaine de définition. Il existe alors un $S$-morphisme et un seul $U_0\to Y$
appartenant à la classe $f$.
\end{proposition}

Comme pour tout morphisme $U\to Y$ appartenant à la classe $f$, on a
nécessairement $U\subset U_0$, il est clair que la proposition sera
conséquence du

\begin{lemma}[7.2.2.1]
\label{I.7.2.2.1-fr}
Sous les hypothèses de (\hyperref[I.7.2.2-fr]{7.2.2}), soient $U_1$, $U_2$
deux ouverts partout denses de $X$, $f_i:U_i\to Y$ ($i=1,2$) deux
$S$-morphismes tels qu'il existe un ouvert $V\subset U_1\cap U_2$, dense dans
$X$ et dans lequel $f_1$ et $f_2$ coïncident. Alors $f_1$ et $f_2$ coïncident
dans $U_1\cap U_2$.
\end{lemma}

On peut évidemment se borner au cas où $X=U_1=U_2$. Comme $X$ (donc $V$)
est réduit, $X$ est le plus petit sous-préschéma fermé de $X$ majorant $V$
(\hyperref[I.5.2.2-fr]{5.2.2}). Soit
$g=(f_1,f_2)_S:X\to Y\times_S Y$~; comme par hypothèse la diagonale
$T=\Delta_Y(Y)$ est un sous-préschéma fermé de $Y\times_S Y$,
$Z=g^{-1}(T)$ est un sous-préschéma fermé de $X$
(\hyperref[I.4.4.1-fr]{4.4.1}). Si $h:V\to Y$ est la restriction commune de
$f_1$ et $f_2$ à $V$, la restriction de $g$ à $V$ est $g'=(h,h)_S$, qui se
factorise en $g'=\Delta_Y\circ h$~; comme $\Delta_Y^{-1}(T)=Y$, on a
$g'^{-1}(T)=V$, et par suite $Z$ est un sous-préschéma fermé de $X$ induisant
$V$, donc majorant $V$, ce qui entraîne $Z=X$. De la relation
$g^{-1}(T)=X$, on déduit (\hyperref[I.4.4.1-fr]{4.4.1}) que $g$ se factorise
en $\Delta_Y\circ f$, où $f$ est un morphisme $X\to Y$, ce qui entraîne par
définition du morphisme diagonal que $f_1=f_2=f$.

Il est clair que le morphisme $U_0\to Y$ défini dans
(\hyperref[I.7.2.2-fr]{7.2.2}) est l'unique morphisme de la classe $f$ qui
\emph{ne puisse être prolongé} à un morphisme d'une partie ouverte de $X$
contenant strictement $U_0$. \emph{Sous les hypothèses de
(\hyperref[I.7.2.2-fr]{7.2.2})}, on peut donc \emph{identifier} les
applications rationnelles de $X$ dans $Y$ aux morphismes \emph{non
prolongeables} (à des ouverts strictement plus grands) d'ouverts partout
denses de $X$ dans $Y$. Avec cette identification, la prop.
(\hyperref[I.7.2.2-fr]{7.2.2}) entraîne~:

\begin{corollary}[7.2.3]
\label{I.7.2.3-fr}
Les hypothèses sur $X$ et $Y$ étant celles de
(\hyperref[I.7.2.2-fr]{7.2.2}), soit $U$ un ouvert partout dense de $X$. Il
existe une correspondance biunivoque canonique entre les $S$-morphismes de
$U$ dans $Y$ et les $S$-applications rationnelles de $X$ dans $Y$ définies
en tous les points de $U$.
\end{corollary}

En vertu de (\hyperref[I.7.2.2-fr]{7.2.2}), pour tout $S$-morphisme $f$ de
$U$ dans $Y$, il existe en effet une $S$-application rationnelle et une seule
$\overline f$ de $X$ dans $Y$ qui prolonge $f$.

\begin{corollary}[7.2.4]
\label{I.7.2.4-fr}
Soient $S$ un schéma, $X$ un $S$-préschéma réduit, $Y$ un $S$-schéma,
$f:U\to Y$ un $S$-morphisme d'un ouvert dense $U$ de $X$ dans $Y$. Si
$\overline f$ est la $\mathbf{Z}$-application rationnelle de $X$ dans $Y$ qui
prolonge $f$, $\overline f$ est un $S$-morphisme (et est par suite la
$S$-application rationnelle de $X$ dans $Y$ prolongeant $f$).
\end{corollary}

En effet, si $\varphi:X\to S$, $\psi:Y\to S$ sont les morphismes structuraux,
$U_0$ le domaine de définition de $\overline f$, $j$ l'injection $U_0\to X$,
il suffit de prouver que $\psi\circ\overline f=\varphi\circ j$, ce qui résulte
aussitôt de (\hyperref[I.7.2.2.1-fr]{7.2.2.1}), puisque $f$ est un
$S$-morphisme.

\begin{corollary}[7.2.5]
\label{I.7.2.5-fr}
Soient $X$, $Y$ deux $S$-préschémas~; on suppose $X$ réduit, $X$ et $Y$
séparés sur $S$. Soient $p:Y\to X$ un $S$-morphisme (faisant de $Y$ un
$X$-préschéma), $U$ un ouvert partout dense de $X$, $f$ une $U$-section de
$Y$~; alors l'application rationnelle $\overline f$ de $X$ dans $Y$
prolongeant $f$ est une $X$-section rationnelle de $Y$.
\end{corollary}

